\chapter*{Abstract}
\addcontentsline{toc}{chapter}{Abstract}

Physically based rendering has matured into a reliable foundation for
film production, game engines, and scientific visualization.
Yet the hardware that accelerates rendering -- modern ray tracing GPUs --
remains entirely proprietary: closed ISA, closed microarchitecture,
closed silicon.
This makes independent verification, teaching, and further research
difficult or impossible.

This thesis presents \gonzales{} and \borg{} -- two complementary
open-source systems that together form the first fully open hardware
path tracing architecture backed by actual silicon:

\medskip
\begin{description}
  \item[\gonzales{}] is a production-quality path tracer written in Swift,
    capable of rendering Disney's \emph{Moana} island scene
    (130\,GB of input data) -- to our knowledge the only single-author
    renderer able to do so.
    It serves as the software reference implementation and benchmark driver.

  \item[\borg{}] is a Tile-Based Deferred Rendering (TBDR) GPU
    implemented in Chisel, synthesized on iCE40 FPGAs and taped out in
    GlobalFoundries 180\,nm silicon via the Tiny Tapeout programme.
    It is, to our knowledge, the first open-source GPU
    for which a physical tapeout has been completed.
\end{description}

\medskip
We extend \borg{}'s instruction set architecture with ray-scene
intersection primitives and demonstrate path tracing workloads
running on \borg{} hardware.
The resulting system provides a fully reproducible, peer-reviewable
hardware platform for rendering research -- closing the gap that has
existed since the earliest ray tracing hardware proposals.

We evaluate the system against CPU-based rendering (Intel Embree) and
commercial GPU ray tracing (NVIDIA RTX), and discuss architectural
trade-offs for embedded and FPGA-class devices.
